\documentclass[11pt]{article}
\usepackage[a4paper,margin=25mm]{geometry}
\usepackage{amsmath,amssymb,amsthm}
\usepackage[hidelinks]{hyperref}
\emergencystretch=2em
\newtheorem{theorem}{Theorem}
\title{An Explicit Finite-Time Three-Vortex Collapse:\\A Disproof of Conjecture 00000005268}
\author{GitHub: gaochengzhecpu}
\date{}
\begin{document}
\maketitle
\begin{abstract}
Three nonzero, mixed-sign circulations admit an explicit solution of the
planar point-vortex equations whose vortices remain pairwise distinct
until time $1/2$ and converge to a common point at that time. The
construction uses rational initial coordinates and a logarithmically
rotating, square-root shrinking similarity factor. Lean verifies the
actual complex-valued differential equation, including its conventional
$1/(2\pi)$ factor, and the full one-sided collision limits.
\end{abstract}

\paragraph{Source and reading.}
Both source languages assert that finite-time collision of three point
vortices is impossible for general strength triples. We interpret
``general'' as imposing no same-sign restriction. A solution that is
collision-free on $[0,T)$ and has a common finite limit as $t\uparrow T$
is a finite-time total collision. The equation need not remain defined
at the collision itself. One such solution disproves the universal
exclusion, regardless of the proposed additional first integral.
This does not assert that collisions occur for same-sign strengths.

\paragraph{Equation and circulations.}
Identify the plane with $\mathbb C$ and use the usual convention
\begin{equation}\label{eq:physical}
 \dot z_i=\frac{1}{2\pi}\sum_{j\ne i}
 K_j\frac{\mathrm i(z_i-z_j)}{|z_i-z_j|^2},\qquad i=0,1,2.
\end{equation}
Set $K_j=2\pi\gamma_j$, where
\[
 (\gamma_0,\gamma_1,\gamma_2)
 =\left(-\frac{1073}{84},-\frac{1073}{140},\frac{1073}{224}\right).
\]
All three physical circulations are nonzero. With this choice,
\eqref{eq:physical} is precisely $\dot z=F(z)$, where
\[
 F_i(z)=\sum_{j\ne i}\gamma_j
 \frac{\mathrm i(z_i-z_j)}{|z_i-z_j|^2}.
\]
This passage is a change of circulation units, not a change in the
physical equation or the meaning of collision.

\paragraph{A rational similarity configuration.}
Define
\[
 c=\frac38-\frac{15}{56}\mathrm i,\qquad
 (w_0,w_1,w_2)=
 \left(-\frac38+\frac{15}{56}\mathrm i,
 \frac58+\frac{15}{56}\mathrm i,\frac87\mathrm i\right),
 \qquad \lambda=-1-\frac{238}{15}\mathrm i.
\]
The squared pairwise distances of $w$ are respectively
$1,29/32,37/32$, so the points are distinct. Direct rational arithmetic
in the displayed definition of $F$ gives
\[
\begin{array}{c|c|c}
i&F_i(w)&\lambda w_i\\\hline
0&37/8+(1591/280)\mathrm i&37/8+(1591/280)\mathrm i\\
1&29/8-(1711/168)\mathrm i&29/8-(1711/168)\mathrm i\\
2&272/15-(8/7)\mathrm i&272/15-(8/7)\mathrm i
\end{array}
\]
Thus $F_i(w)=\lambda w_i$. For any $u\ne0$ and any $c\in\mathbb C$,
the actual vortex field satisfies the similarity identity
\begin{equation}\label{eq:similarity}
 F_i(c+uw)=\frac{u}{|u|^2}F_i(w).
\end{equation}
Indeed, each difference acquires a factor $u$, and its squared modulus
acquires the factor $|u|^2$.

\begin{theorem}
The following solution of \eqref{eq:physical} undergoes a total collision
at time $T=1/2$:
\[
 q=\frac12+\frac{119}{15}\mathrm i,\qquad
 u(t)=\exp\bigl(q\log(1-2t)\bigr),\qquad
 z_i(t)=c+u(t)w_i\quad(0\le t<1/2).
\]
Its initial positions are $0$, $1$, and $3/8+(7/8)\mathrm i$.
\end{theorem}
\begin{proof}
For $t<1/2$, let $s=1-2t>0$. Complex exponentials never vanish, and
\[
 |u(t)|^2=\exp(2\operatorname{Re}(q)\log s)=s,
 \qquad
 u'(t)=\frac{-2q}{s}u(t)=\frac{\lambda}{s}u(t).
\]
Consequently the three positions remain distinct. Differentiation and
\eqref{eq:similarity} give, for each $i$,
\[
 \dot z_i(t)=\frac{\lambda u(t)}s w_i
 =\frac{u(t)}{|u(t)|^2}F_i(w)=F_i(z(t)),
\]
which is exactly \eqref{eq:physical} for $K_j=2\pi\gamma_j$.
Since $u(0)=1$, adding $c$ to the listed $w_i$ gives the initial
positions in the theorem. Finally $|u(t)|=\sqrt{1-2t}\to0$ as
$t\uparrow1/2$. Hence $z_i(t)\to c$ for all three vortices.
\end{proof}

\paragraph{Formal verification and scope.}
The Lean file defines the vector field by an actual sum over
\texttt{Fin 3}, using Mathlib complex numbers and their squared norm.
It proves the rational field identity, similarity covariance,
nonzero circulations, exact initial positions, injectivity of the
configuration before collision, and the actual derivative statements
\texttt{HasDerivAt}. The collision is a Mathlib filter limit from the
left. The final theorem negates a universally quantified proposition
excluding any such solution for all nonzero circulation triples,
all positive finite times, all trajectories, and all collision points.
The conventional $1/(2\pi)$ equation and its circulation rescaling
are also proved in Lean. No numerical simulation, unverified ODE
existence assertion, or endpoint substitution replaces the trajectory.

The separate Hamiltonian formulation and general classification of
collapsing triples are not formalized or needed here. The paper's
direct proof and the Lean theorem concern this explicit solution.
The axiom audit permits only \texttt{propext},
\texttt{Classical.choice}, and \texttt{Quot.sound}.

\paragraph{Source and reproduction.}
The exact bilingual source is preserved as \texttt{SOURCE.md}; the
frozen upstream commit and source digest are recorded in
\texttt{verification/SOURCE\_PROVENANCE.md}.
The formal proof is in \texttt{lean/Main.lean}, using Lean 4.19.0 and
the pinned Mathlib dependency recorded in the project manifest.
From the Lean project directory, run
\begin{center}\texttt{lake env lean -DwarningAsError=true Main.lean}\end{center}
The accompanying strict log contains the actual compiler result and
axiom output. Development was by a delegated AI agent with a
self-review; the coordinating agent reviews it separately. No external
independent review is claimed.
\end{document}
